\documentclass[11pt]{article}
\usepackage[margin=1.1in]{geometry}
\usepackage{amsmath,amssymb,amsthm}
\usepackage[colorlinks,linkcolor=blue,citecolor=blue]{hyperref}

\newtheorem{theorem}{Theorem}
\newtheorem{lemma}[theorem]{Lemma}
\newtheorem{proposition}[theorem]{Proposition}
\newtheorem{corollary}[theorem]{Corollary}
\theoremstyle{definition}
\newtheorem{remark}[theorem]{Remark}
\newtheorem{conjecture}[theorem]{Conjecture}

\newcommand{\R}{\mathbb{R}}
\newcommand{\Sig}{\Sigma}
\DeclareMathOperator{\arctanop}{arctan}

\title{Separation of the two niches of Romik's ambidextrous sofa,\\
and a closed form for Romik's constant $f_1$}
\author{Vico Bonfioli\\ \small\texttt{vicobonfioli@gmail.com}}
\date{July 2026}

\begin{document}
\maketitle

\begin{abstract}
Baek's proof that Gerver's sofa is optimal represents a moving sofa as a convex
cap minus a single niche, and obtains global optimality from the concavity of a
quadratic upper bound. Transferring that architecture to the ambidextrous moving
sofa problem, which remains open, meets an immediate obstruction: the
ambidextrous sofa is a convex cap minus a \emph{union of two} niches, and the
inclusion--exclusion term enters the upper bound with the sign opposite to the
one concavity requires. We show the obstruction is vacuous. The two niches of
Romik's candidate $\Sig$ are disjoint, because every niche point lies at or below
the apex of its own wedge and the apex path stays strictly below the axis of the
reflection exchanging the two niches. The relevant constant is
$M=\max_t c_y(t)=\tfrac12-\bigl(\sqrt2-f_1\sqrt{4-2\sqrt2}\bigr)$, and the
condition $M<\tfrac12$ is equivalent to $f_1^2<(2+\sqrt2)/2$, which holds with a
margin of $0.2600\ldots$. Along the way we give a closed form for Romik's
constant $f_1$, which appears in the literature only as a decimal.
\end{abstract}

\section{Setting}

Write $L$ for the closed L-shaped hallway of unit width with its corner at the
origin, $\mu_t=(\cos t,\sin t)$, $\nu_t=(-\sin t,\cos t)$, and for a rotation
path $c:[0,\pi/2]\to\R^2$ let
\[
  C_t=\{p:\langle p-c(t),\mu_t\rangle\le 1\}\cap\{p:\langle p-c(t),\nu_t\rangle\le 1\},
  \qquad
  Q_t=\{p:\langle p-c(t),\mu_t\rangle<0\}\cap\{p:\langle p-c(t),\nu_t\rangle<0\},
\]
so that the hallway in position $t$ is $H_t=C_t\setminus Q_t$. The one-sided
sofa is $S=\bigcap_t H_t$, and writing $C=\bigcap_t C_t$ and
$U=\bigcup_t Q_t$ we have $S=C\setminus U$ with $C$ convex.

Let $\rho(x,y)=(x,1-y)$ be the reflection in the line $y=\tfrac12$. Romik's
ambidextrous sofa is
\begin{equation}\label{eq:sigma}
  \Sig=S\cap\rho(S)=C_2\setminus\bigl(U\cup\rho U\bigr),
  \qquad C_2:=C\cap\rho C \text{ convex},
\end{equation}
with $|\Sig|=A_R^\ast=1.6449552184\ldots$, and $c$ is Romik's piecewise path with
breakpoints at $\beta$ and $\pi/2-\beta$, where
\begin{equation}\label{eq:beta}
  \beta=\arctanop\Bigl(\tfrac12\bigl(\sqrt[3]{\sqrt2+1}-\sqrt[3]{\sqrt2-1}\bigr)\Bigr)
      =0.2896538208\ldots
\end{equation}

Baek's argument for the one-corner problem builds an overestimating region whose
area is a quadratic functional on a convex domain of convex bodies, and proves it
concave by exhibiting it as a linear functional minus Mamikon regions, whose areas
are convex quadratics. Concavity reduces global optimality to a first-order
condition. The construction is organised around $S=K\setminus\mathcal N(K)$: one
cap, one niche.

By \eqref{eq:sigma} the ambidextrous analogue has \emph{two} niches, and
\begin{equation}\label{eq:ie}
  |U\cup\rho U|=|U|+|\rho U|-|U\cap\rho U|.
\end{equation}
The last term appears in the upper bound with a $+$ sign. Since concavity is
obtained by \emph{subtracting} convex quadratics, an added overlap term is not
covered by the mechanism. The purpose of this note is Theorem~\ref{thm:main}:
the term is zero.

\section{The niches are disjoint}

\begin{lemma}[Niche ceiling]\label{lem:ceiling}
For every $t\in[0,\pi/2]$ we have $Q_t\subseteq\{(x,y):y\le c_y(t)\}$.
\end{lemma}

\begin{proof}
$Q_t$ is the open convex cone with apex $c(t)$ spanned by $-\mu_t$ and $-\nu_t$,
so $q\in Q_t$ means $q-c(t)=-a\mu_t-b\nu_t$ with $a,b>0$. For $t\in[0,\pi/2]$
both $\sin t\ge0$ and $\cos t\ge0$, hence
\[
  q_y-c_y(t)=-a\sin t-b\cos t\le 0. \qedhere
\]
\end{proof}

\begin{corollary}\label{cor:sep}
Put $M:=\max_{t\in[0,\pi/2]}c_y(t)$. Then $U\subseteq\{y\le M\}$ and
$\rho U\subseteq\{y\ge 1-M\}$. If $M<\tfrac12$ then $U\cap\rho U=\emptyset$.
\end{corollary}

\begin{proof}
The first inclusion is Lemma~\ref{lem:ceiling} applied to each $t$; the second is
its image under $\rho$. If $M<\tfrac12$ then $1-M>\tfrac12>M$, so the two
half-planes are disjoint.
\end{proof}

It remains to evaluate $M$ for Romik's path. On the middle piece
$[\beta,\pi/2-\beta]$ Romik's path is
\begin{equation}\label{eq:x6}
  c(t)=R(t)v(t)+\kappa,\qquad
  v(t)=\begin{pmatrix} f_1\cos\frac t2+f_2\sin\frac t2-1\\[2pt]
                      -f_2\cos\frac t2+f_1\sin\frac t2-1\end{pmatrix},
\end{equation}
with $R(t)$ the rotation by $t$, and the two facts we use are
\begin{equation}\label{eq:kf}
  \kappa_y=\tfrac12 \quad\text{exactly},
  \qquad f_2=(1-\sqrt2)\,f_1 .
\end{equation}

\begin{proposition}\label{prop:cy}
For $t\in[\beta,\pi/2-\beta]$,
\begin{equation}\label{eq:cy}
  c_y(t)-\tfrac12
   = f_1\sin\tfrac{3t}{2}-f_2\cos\tfrac{3t}{2}-(\sin t+\cos t)
   = f_1\sqrt{4-2\sqrt2}\,\sin\!\Bigl(\tfrac{3t}{2}+\tfrac\pi8\Bigr)-\sqrt2\,\sin\!\Bigl(t+\tfrac\pi4\Bigr).
\end{equation}
In particular $c_y$ is symmetric about $t=\pi/4$ on this interval, and
\begin{equation}\label{eq:M}
  M=c_y(\pi/4)=\tfrac12-\Bigl(\sqrt2-f_1\sqrt{4-2\sqrt2}\Bigr).
\end{equation}
\end{proposition}

\begin{proof}
By \eqref{eq:x6} and $\kappa_y=\tfrac12$, $c_y(t)-\tfrac12=\sin t\,v_x+\cos t\,v_y$.
Write $S=\sin\frac t2$, $C=\cos\frac t2$, so $\sin t=2SC$ and $\cos t=C^2-S^2$.
Expanding and collecting the coefficients of $f_1$ and $f_2$ gives
\[
  \sin t\,v_x+\cos t\,v_y
  = f_1 S(3C^2-S^2)+f_2 C(3S^2-C^2)-(\sin t+\cos t),
\]
and the triple-angle identities $S(3C^2-S^2)=\sin\frac{3t}{2}$,
$C(3S^2-C^2)=-\cos\frac{3t}{2}$ give the first equality in \eqref{eq:cy}. For the
second, substitute $f_2=(1-\sqrt2)f_1$ from \eqref{eq:kf}: the bracket becomes
$\sin\frac{3t}{2}+(\sqrt2-1)\cos\frac{3t}{2}$, of amplitude
$\sqrt{1+(\sqrt2-1)^2}=\sqrt{4-2\sqrt2}$ and phase $\arctanop(\sqrt2-1)=\pi/8$.

Both sinusoids in \eqref{eq:cy} are invariant under $t\mapsto\pi/2-t$: indeed
$\sin(\tfrac{3}{2}(\pi/2-t)+\pi/8)=\sin(7\pi/8-\tfrac{3t}{2})=\sin(\pi/8+\tfrac{3t}{2})$
and $\sin(3\pi/4-t)=\sin(\pi/4+t)$. Hence $c_y$ is symmetric about $\pi/4$, where
both sinusoids attain the value $1$ simultaneously; \eqref{eq:M} follows. \end{proof}

That $\pi/4$ maximises $c_y$ over all of $[0,\pi/2]$, and not merely over the
middle piece, follows from the corresponding closed form on the outer pieces.

\begin{proposition}\label{prop:outer}
For $t\in[0,\beta]$,
\begin{equation}\label{eq:cyouter}
  c_y(t)=\tfrac12+a_1\sin 2t-\sin t-\tfrac12\cos t,
\end{equation}
and $c_y$ is strictly increasing there, so $\max_{[0,\beta]}c_y=c_y(\beta)$. By
the symmetry $c_y(t)=c_y(\pi/2-t)$ the same bound holds on $[\pi/2-\beta,\pi/2]$.
Consequently $M=c_y(\pi/4)$ as claimed in \eqref{eq:M}.
\end{proposition}

\begin{proof}
On $[0,\beta)$ Romik's path is $c(t)=R(t)v_1(t)+\kappa^{(1)}$ with
$v_1(t)=(a_1\cos t-1,\,a_1\sin t-\tfrac12)$ and $\kappa^{(1)}_y=\tfrac12$, so
$c_y(t)-\tfrac12=\sin t\,v_{1x}+\cos t\,v_{1y}=2a_1\sin t\cos t-\sin t-\tfrac12\cos t$,
which is \eqref{eq:cyouter}. Differentiating,
\[
  c_y'(t)=2a_1\cos 2t-\cos t+\tfrac12\sin t
        \ \ge\ 2a_1\cos 2\beta-1\ =\ 0.4650\ldots>0
        \qquad (t\in[0,\beta]),
\]
using $\cos 2t\ge\cos 2\beta$ and $\cos t-\tfrac12\sin t\le 1$ for $t\ge0$.
Hence $c_y$ is strictly increasing on $[0,\beta]$, with
$c_y(\beta)=0.2143795161\ldots$, which is below
$c_y(\pi/4)=0.3878381292\ldots$.
\end{proof}

\begin{remark}
Evaluating \eqref{eq:cyouter} and \eqref{eq:cy} at $t=\beta$ gives the same value
to $4.4\cdot 10^{-31}$ at $30$-digit precision. Since the two expressions were
derived from different pieces of Romik's path, this is an independent check on
both, and on \eqref{eq:f1}.
\end{remark}

\section{A closed form for $f_1$}

Romik's constant $f_1$ is recorded in the literature as the decimal
$1.202938908156911389$. The following determines it in closed form. Let
\begin{equation}\label{eq:a1}
  a_1=\tfrac14\sqrt{\,4+\sqrt[3]{71+8\sqrt2}+\sqrt[3]{71-8\sqrt2}\,}
     =0.8752873624\ldots
\end{equation}

\begin{proposition}\label{prop:f1}
\begin{equation}\label{eq:f1}
  f_1=\frac{\tfrac43\,a_1\cos\beta}{\cos\frac\beta2+(1-\sqrt2)\sin\frac\beta2}.
\end{equation}
\end{proposition}

\begin{proof}
On $[0,\beta)$ Romik's path is $c(t)=R(t)v_1(t)+\kappa^{(1)}$ with
$v_1(t)=(a_1\cos t-1,\;a_1\sin t-\tfrac12)$ and
$\kappa^{(1)}=(1-a_1,\tfrac12)$; on $[\beta,\pi/2-\beta]$ it is \eqref{eq:x6}
with $\kappa=(1-\tfrac43 a_1,\tfrac12)$. Both $\kappa$'s have second coordinate
$\tfrac12$ and their first coordinates differ by exactly $\tfrac13 a_1$.
Continuity at $t=\beta$ therefore gives
$R(\beta)\bigl(v_1(\beta)-v(\beta)\bigr)=(-\tfrac13a_1,0)$, so
\[
  v_1(\beta)-v(\beta)=R(-\beta)\bigl(-\tfrac13a_1,0\bigr)
     =\bigl(-\tfrac13a_1\cos\beta,\;\tfrac13a_1\sin\beta\bigr).
\]
Reading the first coordinate and using $f_2=(1-\sqrt2)f_1$,
\[
  \bigl(a_1\cos\beta-1\bigr)-\Bigl(f_1\bigl(\cos\tfrac\beta2+(1-\sqrt2)\sin\tfrac\beta2\bigr)-1\Bigr)
   =-\tfrac13a_1\cos\beta,
\]
which rearranges to \eqref{eq:f1}.
\end{proof}

\begin{remark}
Evaluating \eqref{eq:f1} at $50$ digits gives
$f_1=1.202938908156911389070223\ldots$, so the tabulated decimal is its rounding.
As an independent check, the \emph{second} coordinate of the junction relation,
which was not used in the derivation, is satisfied to $1.7\cdot10^{-51}$ at the
same precision.
\end{remark}

\section{Main theorem}

\begin{theorem}\label{thm:main}
For Romik's ambidextrous sofa, $U\cap\rho U=\emptyset$, and consequently
\[
  |\Sig|=|C_2|-|U|-|\rho U|=|C_2|-2|U| .
\]
\end{theorem}

\begin{proof}
By Proposition~\ref{prop:cy}, $M<\tfrac12$ is equivalent to
$f_1\sqrt{4-2\sqrt2}<\sqrt2$, i.e. to
\begin{equation}\label{eq:ineq}
  f_1^2<\frac{2+\sqrt2}{2}.
\end{equation}
With $f_1$ given by \eqref{eq:f1} in terms of the algebraic numbers
\eqref{eq:beta} and \eqref{eq:a1}, both sides of \eqref{eq:ineq} are explicit
algebraic numbers, and
\[
  f_1^2=1.4470620167577420940\ldots,
  \qquad \frac{2+\sqrt2}{2}=1.7071067811865475244\ldots,
\]
so \eqref{eq:ineq} holds with margin $0.2600447644\ldots$. Corollary~\ref{cor:sep}
gives $U\cap\rho U=\emptyset$, and then \eqref{eq:ie} together with
\eqref{eq:sigma}. Finally $|\rho U|=|U|$ because $\rho$ is an isometry.
\end{proof}

\begin{corollary}
$M=0.3878381292441942964\ldots$, the two niches are separated by the horizontal
band of width $1-2M=0.2243237415116114\ldots$ about $y=\tfrac12$, and the
ambidextrous upper bound requires only \emph{one} niche to be modelled: one core
and two tails, the same shape as in the one-corner problem.
\end{corollary}

\section{The quotient reduction and a conditional statement}

Theorem~\ref{thm:main} has a consequence that changes the shape of the
ambidextrous problem. Since $U\subseteq\{y\le M\}$ with $M<\tfrac12$, the whole
lower niche lies below the symmetry axis, so intersecting with the lower
half-strip deletes $\rho U$ outright.

\begin{corollary}[Quotient reduction]\label{cor:quot}
Put $K^-:=C_2\cap\{y\le\tfrac12\}$, a convex body. Then
\[
  \Sig\cap\{y\le\tfrac12\}=K^-\setminus U,
  \qquad
  |\Sig|=2\bigl(|K^-|-|U\cap K^-|\bigr).
\]
\end{corollary}

That is a convex cap minus \emph{one} niche: the shape Baek's argument is built
on. Numerically, at $n=1441$ hallways, $|\Sig|=1.645059395$ and
$|\Sig\cap\{y\le\frac12\}|=0.822529698$, whose double reproduces $|\Sig|$ to
$3.8\cdot10^{-15}$.

Two cautions must be recorded, because they bound what
Corollary~\ref{cor:quot} delivers.

First, $M<\tfrac12$ is a property of Romik's trajectory, not of an arbitrary
competitor. A competitor with $M\ge\tfrac12$ has overlapping niches, and by
\eqref{eq:ie} overlap \emph{increases} $|\Sig|$; so it cannot be assumed away on
extremal grounds.

Second --- and this is what makes the reduction usable --- connectedness does
appear to exclude it, but not for the reason one first tries. A \emph{single}
pair $Q_t\cup\rho Q_t$ removes a full vertical segment, and hence disconnects any
body meeting both sides of it, exactly when the apex height exceeds a threshold
$h^\ast(t)$; and computation gives $h^\ast(t)>\tfrac12$ throughout $(0,\pi/2)$,
with $h^\ast(t)\to\tfrac12$ as $t\to0$, $h^\ast(\pi/4)=0.5100$ and
$h^\ast(1.52)=0.6967$. So no single-$t$ argument forces $M\le\tfrac12$; all it
gives is the pointwise constraint $c_y(t)\le h^\ast(t)$.

The full family does force it. Deforming Romik's path upward near its maximum,
$c_y\mapsto c_y+\delta g$ with $g$ vanishing at both ends, and computing
$\Sig=S\cap\rho S$, one finds that $\Sig$ is a single connected region for
$M<\tfrac12$ and splits into several components for $M>\tfrac12$, with the
transition at $M=\tfrac12$ to six decimal places: at $M=0.4999981$ the body is
connected, at $M=0.5000081$ it has four components. The threshold is unchanged
under variation of the bump's location and width (peaks at $0.5$, $0.6$, $1.0$,
$\pi/4$ and widths $0.4$ to $0.9$ all give connected just below $\tfrac12$ and
split just above).

\begin{conjecture}\label{conj:conn}
Every connected ambidextrous moving sofa satisfies $\max_t c_y(t)\le\tfrac12$.
\end{conjecture}

The mechanism is visible at the threshold. When $c_y(t_0)=\tfrac12$ the apex lies
on the symmetry axis, so $\rho c(t_0)=c(t_0)$ and the two wedges share an apex.
Their angular spans are $[t_0+\pi,\,t_0+\tfrac{3\pi}{2}]$ and
$[\tfrac\pi2-t_0,\,\pi-t_0]$, which leave an angular gap of width $2t_0$; that is
why the pair alone does not separate. For $c_y(t_0)>\tfrac12$ the wedges of
nearby $t$, together with their $\rho$-images, close the gap. A proof of
Conjecture~\ref{conj:conn} along these lines would discharge the hypothesis below.

\begin{proposition}[Conditional form]\label{prop:cond}
Let $S$ be an ambidextrous moving sofa with rotation path $c$ satisfying
$\max_t c_y(t)\le\tfrac12$. Then its two niches are disjoint and
$|S|=2(|K^-|-|U\cap K^-|)$ with $K^-$ convex and $U$ a single niche. In
particular the area functional of the ambidextrous problem restricted to this
class is of the cap-minus-one-niche form to which the concavity method applies.
Granting Conjecture~\ref{conj:conn}, the hypothesis is automatic and the
conclusion holds for every ambidextrous moving sofa.
\end{proposition}

We emphasise what is \emph{not} claimed: no optimality assertion follows from
Proposition~\ref{prop:cond} alone. Baek's route additionally requires an
existence and regularity step (a maximum monotone sofa of rotation angle $\pi/2$
satisfying an injectivity condition), and his balancing argument for that step
does not transfer, for a reason worth recording. His argument turns on a
coincidence: the derivative of the area functional under pushing an edge is
$\sigma(t)-\tau(t)$, and the same combination is forced to sum to zero because
the cap boundary and the niche polyline join the \emph{same} pair of endpoints.
In the quotient, $\partial K^-$ acquires the horizontal cut at $y=\tfrac12$,
where the polyline contributes nothing, and the identity becomes
$\sum_{t\ne\mathrm{cut}}(\sigma-\tau)(v_t\cdot u_0)=\sigma_{\mathrm{cut}}>0$,
which is consistent with $\sigma\le\tau$ throughout. One obtains an inequality
where Baek obtains an equality, and it is the equality that his argument
consumes.

\section{Formalisation}

Lemma~\ref{lem:ceiling}, Corollary~\ref{cor:sep} and the collapse of
\eqref{eq:ie} are machine-checked in Lean~4 without Mathlib as
\verb|niche_below_apex|, \verb|niche_disjoint| and
\verb|union_area_of_disjoint|. The identity underlying $\Sigma$'s middle-phase
equation is also formalised: writing $D$ for the formal derivative in the
half-angle, $D(Dv)=-(v+1)$ for every arc of constant term $-1$, which is
$4v''=-(v+(1,1))$ for the bracket of \eqref{eq:x6}
(\verb|Trig.sol6_bracket|, proved by \verb|rfl| and depending on no axioms). \verb|lake build| is clean with no \verb|sorry|,
and \verb|#print axioms| reports nothing beyond \verb|propext| and
\verb|Quot.sound|. Trigonometric quantities are carried as formal symbols with
the sign hypotheses $\sin t\ge0$, $\cos t\ge0$, which is all
Lemma~\ref{lem:ceiling} uses.

\begin{thebibliography}{9}
\bibitem{Baek} J.~Baek, \emph{Optimality of Gerver's sofa},
  arXiv:2411.19826, 2024.
\bibitem{Gerver} J.~L.~Gerver, \emph{On moving a sofa around a corner},
  Geom.\ Dedicata \textbf{42} (1992), 267--283.
\bibitem{Romik} D.~Romik, \emph{Differential equations and exact solutions in the
  moving sofa problem}, Exp.\ Math.\ \textbf{27} (2018), 316--330.
\end{thebibliography}

\end{document}
